\chapter{Method A --- Averaged mode}
\label{sec:avg}

\begin{center}
\small
\begin{tabular}{@{}ll@{}}
\toprule
Notebook cell & Step 4A (analysis: Step 5A) \\
Code path & \co{prog.acquire(progress=False)} \\
One call returns & the \emph{mean} over \co{reps} reps: one sample \\
Measured rate & \textbf{86.8--88.8\,Hz}, at every $\tau$ and rep count tested \\
Rate limited by & the host call ($\sim$11.4\,ms), \emph{not} the readout window \\
Duty cycle & 20.9\% at 10 reps, 48.3\% at 23 reps, 85.7\% at $\tau$=213/23 reps \\
Measured floor & 274--501\,\etaunit\ (varies 1.8$\times$ between runs) \\
Open defects & none \\
\bottomrule
\end{tabular}
\end{center}

\section{Mechanics}

The program runs \co{reps} repetitions. \co{analyze_results()} averages over the
rep axis (Eq.~\ref{eq:divisions} divides by \co{reps}) and returns one number per
parked frequency. The host then converts that pair into one $\Delta f$ and one CSV
row, and immediately calls \co{acquire()} again.

\paragraph{The cost of this design.} The rep axis is \emph{discarded}. Those
\co{reps} repetitions were already independent time samples spaced by
$t_\text{rep}$; averaging them away throws that time resolution out. Methods B and
C exist to keep it.

\section{Timing budget, component by component}
\label{sec:avgbudget}

Three operating points, all measured. Read the \textbf{Adds?} column first: it is
what distinguishes this mode from the other two.

\subsection*{Operating point 1: $\tau=\SI{120}{\micro\second}$, 10 reps --- what actually ran on 2026-08-14}

\begin{center}
\small
\begin{tabular}{@{}>{\raggedright\arraybackslash}p{4.7cm}rrc>{\raggedright\arraybackslash}p{4.2cm}@{}}
\toprule
\hd{Component} & \hd{Time} & \hd{Share} & \hd{Adds?} & \hd{How obtained} \\
\midrule
ADC integration ($10\times\SI{240}{\micro\second}$)
  & \SI{2.400}{\milli\second} & 20.9\% & \textcolor{bad}{no} & exact, Eq.~\ref{eq:recover}, $D$=368640 \\
\quad of which MW pulse
  & \SI{0}{\milli\second} & 0\% & no & concurrent with the readout \\
\quad of which relax/sync
  & $\sim$\SI{0}{\milli\second} & $\sim$0\% & no & absorbed; measured, \S\ref{sec:asm} \\
\midrule
\co{acquire()} call, fastest observed
  & \SI{10.302}{\milli\second} & 89.8\% & --- & p05 of \co{acq_seconds} \\
\co{acquire()} call, median
  & \textbf{\SI{11.408}{\milli\second}} & \textbf{99.4\%} & \textbf{yes} & median of \co{acq_seconds} \\
\co{acquire()} jitter (p95$-$p05)
  & \SI{1.624}{\milli\second} & 14.2\% & yes & host/network scheduling \\
\midrule
Python outside \co{acquire()}
  & \SI{0.067}{\milli\second} & 0.58\% & yes & period $-$ \co{acq_seconds} \\
CSV flush stall (worst)
  & \SI{215.8}{\milli\second} & --- & yes & 2 events $>3\times$ median period \\
\midrule
\textbf{TOTAL per sample} & \textbf{\SI{11.475}{\milli\second}} & 100\% & & $\rightarrow$ \textbf{87.1\,Hz} \\
\bottomrule
\end{tabular}
\end{center}

\begin{keybox}[Read this table as a maximum, not a sum]
$2.400 + 11.408 = 13.8$\,ms, but the measured period is 11.475\,ms. The two do
not add, because the ADC integration happens \emph{inside} stage 6 of the host
chain (\S\ref{sec:rpc}). The period is
\begin{align*}
\text{period} &= \max\bigl(\text{host chain},\ \text{reps}\times t_\text{rep}\bigr)
   \;+\; \text{Python outside} \\
   &= \max(11.408,\ 2.400) + 0.067 \;=\; \SI{11.475}{\milli\second}.
\end{align*}
\textbf{Only \SI{2.400}{\milli\second} of every \SI{11.475}{\milli\second} is spent
collecting photons.} The remaining 79\% is the host talking to the board.
\end{keybox}

\subsection*{Operating point 2: $\tau=\SI{120}{\micro\second}$, 23 reps}

\begin{center}
\small
\begin{tabular}{@{}lrrl@{}}
\toprule
\hd{Component} & \hd{Time} & \hd{Share} & \hd{Note} \\
\midrule
ADC integration: $23\times\SI{240}{\micro\second}$ & \SI{5.520}{\milli\second} & 48.3\% & $D$=847872, exact \\
\co{acquire()} call, fastest & \SI{10.266}{\milli\second} & 89.7\% & \\
\co{acquire()} call, median & \SI{11.384}{\milli\second} & 99.5\% & \\
\co{acquire()} jitter & \SI{1.819}{\milli\second} & 15.9\% & \\
Python outside & \SI{0.056}{\milli\second} & 0.49\% & \\
CSV flush stalls & \SI{0}{\milli\second} & --- & none in this run \\
\midrule
\textbf{TOTAL per sample} & \textbf{\SI{11.440}{\milli\second}} & 100\% & $\rightarrow$ \textbf{87.4\,Hz} \\
\bottomrule
\end{tabular}
\end{center}

\subsection*{Operating point 3: $\tau=\SI{213}{\micro\second}$, 23 reps --- the 2026-08-06 default}

\begin{center}
\small
\begin{tabular}{@{}lrrl@{}}
\toprule
\hd{Component} & \hd{Time} & \hd{Share} & \hd{Note} \\
\midrule
ADC integration: $23\times\SI{426}{\micro\second}$ & \SI{9.798}{\milli\second} & 85.7\% & $D$=1504982, exact \\
\co{acquire()} call, fastest & \SI{10.218}{\milli\second} & 89.4\% & \\
\co{acquire()} call, median & \SI{11.369}{\milli\second} & 99.5\% & \\
\co{acquire()} jitter & \SI{1.877}{\milli\second} & 16.4\% & \\
Python outside & \SI{0.058}{\milli\second} & 0.51\% & \\
CSV flush stall (worst) & \SI{219.5}{\milli\second} & --- & 4 events \\
\midrule
\textbf{TOTAL per sample} & \textbf{\SI{11.428}{\milli\second}} & 100\% & $\rightarrow$ \textbf{87.5\,Hz} \\
\bottomrule
\end{tabular}
\end{center}

\subsection{The three side by side --- the decisive comparison}

\begin{center}
\small
\begin{tabular}{@{}lrrrr@{}}
\toprule
& \hd{$\tau$=120, 10 reps} & \hd{$\tau$=120, 23 reps} & \hd{$\tau$=213, 23 reps} & \hd{spread} \\
\midrule
ADC integration & \SI{2.400}{\milli\second} & \SI{5.520}{\milli\second} & \SI{9.798}{\milli\second} & \textbf{4.1$\times$} \\
Period          & \SI{11.475}{\milli\second} & \SI{11.440}{\milli\second} & \SI{11.428}{\milli\second} & \textbf{1.004$\times$} \\
Rate            & 87.1\,Hz & 87.4\,Hz & 87.5\,Hz & 1.005$\times$ \\
Duty cycle      & 20.9\% & 48.3\% & 85.7\% & 4.1$\times$ \\
\bottomrule
\end{tabular}
\end{center}

\begin{keybox}[This one row is the whole finding]
The ADC integration time varies by \textbf{4.1$\times$} across these three
operating points. The period varies by \textbf{0.4\%}. There is no readout-window
or rep-count setting that makes this mode faster than $\sim$88\,Hz.
\end{keybox}

\section{The CSV flush, and why it was fixed}

The worst single stall in these runs is \SI{219.5}{\milli\second} --- nineteen
missed samples. The cause was the old save routine calling
\co{pd.DataFrame(rows).to_csv(...)} every 5 seconds, rewriting \emph{every row
collected so far}: $O(N^2)$ over a run. Of the 94 outlier batches across the
2026-08-06 session, \textbf{46 sat within 150\,ms of a 5-second boundary}.

Replaced by \co{twopoint_runner.IncrementalCsvWriter}, which appends only the new
rows (header once). On a 40\,000-row run this is 10.2$\times$ cheaper and the
stall is flat with run length instead of growing.

\section{What to set}

\begin{center}
\small
\begin{tabular}{@{}llp{7.4cm}@{}}
\toprule
\hd{Parameter} & \hd{Recommended} & \hd{Why} \\
\midrule
\co{AVG_REPS_PER_BATCH} & \co{None} & Fills the measured call floor automatically
  via \co{reps_that_fit()}. At $\tau$=120 that is \textbf{42 reps}, against the 10
  actually used --- see \S\ref{sec:freeavg}. \\
\co{AVG_CALIBRATE_FLOOR} & \co{True} & Measures the floor on this rig at startup;
  it already differs by 1.2\,ms between the two sessions. Costs $\sim$0.3\,s. \\
\co{readout_integration_tus} & 120 & Inside the flat sensitivity band
  (\S\ref{sec:window}). Not a rate lever here. \\
\co{AVG_PAIR_ORDER} & \co{"forward"} & \co{"abba"} is free at equal averaging but
  untested (\S\ref{sec:abba}). \\
\co{AVG_SKIP_REFERENCE} & \co{True} & The MW-off reference is contaminated on this
  setup. \\
\co{AVG_SHOW_PLOT} & \co{False} & A live 3-panel redraw costs 20--100\,ms, which
  at an 11\,ms period would dominate the loop. \\
\bottomrule
\end{tabular}
\end{center}

\section{Failure modes and what warns you}

\begin{center}
\small
\begin{tabular}{@{}p{4.6cm}p{5.2cm}p{4.6cm}@{}}
\toprule
\hd{Failure} & \hd{Symptom} & \hd{Guard} \\
\midrule
Resonance drifts off the parked pair &
  Dip depth collapses, $\sigma(\Delta f)$ grows as $1/$depth (Eq.~\ref{eq:noise}) &
  \co{report_calibration_validity()} prints measured vs calibrated depth and
  depth/noise; warns below 50\% or ratio $<3$ \\
60\,Hz mains &
  Appears at 26.8--28.8\,Hz, indistinguishable from signal &
  \co{alias_report()} in Step 5A; \textbf{cannot be removed} (\S\ref{sec:alias}) \\
Averaging does not help &
  $\sigma$ gets \emph{worse} with more reps &
  \Open{} --- see \S\ref{sec:notshot} \\
Peaking transient &
  Per-batch transient from an MW polarisation pulse at each batch boundary &
  \co{pre_init=False} on the live program, enforced by \co{assert}; primed once
  with a throwaway \co{pre_init=True} acquire \\
Stale buffer &
  Would show as a bit-identical repeat &
  Freshness guard; \textbf{zero occurrences} in 20\,365 averaged batches \\
\bottomrule
\end{tabular}
\end{center}

\section{When to use it}

When $\sim$88\,Hz is enough. It is the only one of the three methods with no open
defect, its timestamps are honest, and it has never been observed to replay a
buffer. For drone sensing and MCG, where bandwidth is the point, it is the wrong
choice --- go to Method B or C.
